% ============================================================
% UNIDAD: EMPRENDEDORISMO Y DESARROLLO LOCAL
% ============================================================
\clearpage
\section*{UNIDAD: EMPRENDEDORISMO Y DESARROLLO LOCAL}
\subsection*{Teorías, casos aeronáuticos y territorio}

\begin{recuadro}
\textbf{Objetivo de la unidad:} comprender qué es emprender, qué tipos de emprendimientos
existen, analizar casos reales de la industria aeronáutica argentina y entender cómo los
emprendimientos se relacionan con el desarrollo de un territorio (barrio, municipio, región).
\end{recuadro}

% ------------------------------------------------------------
\section*{1. ¿QUÉ ES EMPRENDER?}

\textbf{Emprender} es identificar una oportunidad o un problema y organizar recursos (personas,
dinero, conocimiento, tiempo) para crear algo nuevo que genere valor: un producto, un servicio,
una empresa, una organización social. El \textbf{emprendedor} asume el riesgo de que el
proyecto no funcione.

\subsubsection*{1.1. Teorías del emprendedorismo}
\begin{itemize}
  \item \textbf{Richard Cantillon (siglo XVIII):} el emprendedor compra a un precio cierto
        para vender a un precio incierto; su rasgo central es \textbf{asumir el riesgo}.
  \item \textbf{Joseph Schumpeter (1911):} el emprendedor es un \textbf{innovador} que
        introduce nuevas combinaciones (nuevos productos, métodos, mercados, materias primas u
        organizaciones). Su idea de \textbf{destrucción creativa}: lo nuevo desplaza a lo viejo
        (el GPS desplazó a muchos equipos de radionavegación).
  \item \textbf{Israel Kirzner:} el emprendedor está \textbf{alerta} a oportunidades que otros
        no ven y aprovecha las diferencias entre lo que el mercado ofrece y lo que necesita.
  \item \textbf{Peter Drucker:} la innovación es una disciplina que se puede aprender y
        sistematizar; no depende solo del talento o la suerte.
  \item \textbf{Enfoques actuales:} el emprendimiento se aprende y se construye probando,
        midiendo y corrigiendo (\textit{lean startup}), y depende también del \textbf{entorno}:
        Estado, universidades, escuelas técnicas, financiamiento y redes.
\end{itemize}

\subsubsection*{1.2. Tipos de emprendimiento}
\begin{center}
\small
\begin{tabular}{|p{0.22\linewidth}|p{0.68\linewidth}|}
\hline
\textbf{Tipo} & \textbf{Características y ejemplos} \\ \hline
Por necesidad & Se emprende porque no hay empleo o el ingreso no alcanza (reparaciones a
domicilio, venta de productos). \\ \hline
Por oportunidad & Se detecta un mercado no atendido y se elige emprender (un taller de
aviónica en una región sin talleres habilitados). \\ \hline
Tecnológico & Basado en conocimiento técnico o científico: desarrollo de equipos, software,
drones, componentes. \\ \hline
Social & Su objetivo principal es resolver un problema social o ambiental; la ganancia es un
medio (reciclado de RAEE, cooperativas de trabajo, inclusión laboral). \\ \hline
Cultural & Producción artística, turismo, patrimonio, museos aeronáuticos, festivales. \\ \hline
Familiar & La propiedad y la gestión están en manos de una familia; el desafío es la
\textbf{sucesión} entre generaciones. \\ \hline
Intraemprendimiento & Un empleado que crea un proyecto nuevo dentro de la empresa donde
trabaja. \\ \hline
\end{tabular}
\end{center}

\subsubsection*{1.3. Actitud emprendedora}
Más que un ``don'', son capacidades que se entrenan: iniciativa, tolerancia a la
incertidumbre y al error, perseverancia, capacidad de aprender rápido, trabajo en equipo,
comunicación y ética. El \textbf{fracaso} es parte del proceso: lo importante es fallar
barato, rápido y aprender.

\paragraph*{Perspectiva crítica.} No todo emprendimiento es una elección: muchas veces es la
respuesta a la falta de empleo, sin protección social ni ingresos estables. Pensar el
emprendedorismo también implica pensar el rol del Estado, el acceso al crédito y las
condiciones de trabajo.

% ------------------------------------------------------------
\section*{2. CASOS DEL SECTOR AERONÁUTICO ARGENTINO}

\subsubsection*{2.1. FAdeA -- Fábrica Argentina de Aviones ``Brig. San Martín'' S.A.}
\begin{itemize}
  \item \textbf{Origen:} nació en Córdoba en 1927 como \textbf{Fábrica Militar de Aviones},
        primera fábrica aeronáutica de América Latina. Allí se diseñaron el Pulqui I y II
        (primeros aviones a reacción de la región), el IA-58 Pucará y el IA-63 Pampa.
  \item \textbf{Privatización y reestatización:} en 1995 fue concesionada a Lockheed Martin y
        en 2009 volvió al Estado como sociedad anónima estatal, con el nombre actual.
  \item \textbf{Actividades:} fabricación del entrenador \textbf{IA-63 Pampa III},
        mantenimiento y modernización de aeronaves militares (C-130 Hércules), servicios para
        terceros y desarrollo de proyectos como el IA-100.
  \item \textbf{Situación reciente:} atraviesa una crisis financiera, con dificultades para
        pagar salarios y una cadena de pymes proveedoras afectada.
\end{itemize}
\textbf{Para analizar:} ¿qué ventajas y qué riesgos tiene una empresa estatal estratégica que
depende casi de un solo cliente (el Estado)? ¿Qué pasa con las pymes proveedoras cuando la
empresa ancla entra en crisis?

\subsubsection*{2.2. Cicaré Helicópteros (Saladillo, Buenos Aires)}
\begin{itemize}
  \item \textbf{El emprendedor:} Augusto ``Pirincho'' Cicaré (1937--2022), inventor
        autodidacta. En 1958 construyó el primer helicóptero de América Latina en su pueblo,
        con recursos mínimos.
  \item \textbf{La empresa:} empresa familiar dedicada a helicópteros ultralivianos y
        simuladores de entrenamiento. Es la única fabricante de helicópteros ultralivianos de
        la región y exporta a varios continentes; abrió una filial en Estados Unidos.
  \item \textbf{Desarrollo local:} prefirió fabricar en Saladillo antes que vender sus
        patentes al exterior; la ciudad fue declarada Capital Provincial del Helicóptero.
  \item \textbf{Situación reciente:} en 2026 entró en concurso preventivo para reestructurar
        sus deudas.
\end{itemize}
\textbf{Para analizar:} ¿qué características del emprendedor aparecen en la historia de
Cicaré? ¿Qué desafíos enfrenta una empresa familiar cuando el fundador ya no está? ¿Por qué
una empresa exitosa técnicamente puede tener problemas financieros?

\subsubsection*{2.3. Redimec (Tandil, Buenos Aires)}
\begin{itemize}
  \item \textbf{Origen:} pyme tecnológica fundada en 1994 por el ingeniero Fabián Oyarbide.
  \item \textbf{Actividades:} modernización de cabinas y aviónica (desde el diagnóstico y el
        diseño hasta la certificación), integración de sistemas, distribución y soporte de
        equipos de aviónica y visión nocturna de fabricantes internacionales.
  \item \textbf{Proyectos:} modernización de la aviónica de los Tucano de la Fuerza Aérea,
        participación en la cabina del IA-100 de FAdeA y en proyectos navales.
\end{itemize}
\textbf{Para analizar:} ¿cómo una pyme del interior llega a trabajar para la defensa y la
aviación? ¿Qué papel juega ser representante de marcas internacionales? ¿Qué perfil técnico
necesita una empresa como esta? (Es un posible empleador para un técnico aviónico).

\subsubsection*{2.4. Papiz Turbine}
Caso a desarrollar en clase con la información que aporte el docente: historia de la empresa,
servicios que presta, clientes, certificaciones y su relación con el entramado aeronáutico
local.

% ------------------------------------------------------------
\subsubsection*{2.5. El Área de Material Quilmes: nuestro propio entorno}
Nuestra escuela funciona dentro del \textbf{Área de Material Quilmes} de la Fuerza Aérea
Argentina. El predio tiene una historia industrial propia: allí funcionaba
\textbf{IMPA} (Industria Metalúrgica y Plástica Argentina), una fábrica que en 1941 abrió un
departamento aeronáutico; de ahí viene el nombre con el que se conoce a la escuela. Hoy en el
Área trabajan organismos y empresas de base tecnológica, lo que convierte al entorno de la
escuela en un ejemplo concreto de articulación entre Estado, industria y conocimiento.

\paragraph*{INVAP S.E.}
\begin{itemize}
  \item \textbf{Origen:} creada en 1976 en San Carlos de Bariloche, a partir del trabajo de
        Conrado Varotto y un grupo de profesionales del Centro Atómico Bariloche, mediante un
        convenio entre la Provincia de Río Negro y la CNEA. Su nombre viene de
        \textit{INVestigación APlicada}. Es una empresa estatal de la Provincia de Río Negro y
        cumplió 50 años en 2026.
  \item \textbf{Qué hace:} trabaja por \textbf{proyectos} tecnológicos complejos, ``llave en
        mano'', en cuatro áreas:
  \begin{itemize}
    \item \textbf{Nuclear:} reactores de investigación y producción de radioisótopos,
          exportados a países como Egipto, Perú, Argelia, Australia, Países Bajos y Arabia
          Saudita.
    \item \textbf{Espacial:} principal contratista de la CONAE (satélites SAC y SAOCOM) y
          constructora de los satélites geoestacionarios de comunicaciones ARSAT-1 y ARSAT-2.
    \item \textbf{Defensa y seguridad:} radares primarios y secundarios para control de tránsito
          aéreo y defensa, radares meteorológicos, sistemas de comando y control.
    \item \textbf{Industria, salud y energía:} equipos de medicina nuclear, sistemas
          industriales, energías renovables.
  \end{itemize}
  \item \textbf{Radares y Fuerza Aérea:} junto con la Fuerza Aérea desarrolló los radares
        secundarios monopulso (RSMA) para el control del tránsito aéreo civil y el
        \textbf{Radar Primario Argentino} 3D de uso militar, dentro del plan SINVICA. Es un caso
        clave para un técnico aviónico: radares, transpondedores y comunicaciones son su
        especialidad.
\end{itemize}
\textbf{Para analizar:} ¿cómo una empresa nacida en un laboratorio nuclear llegó a fabricar
satélites y radares? (``primero el área nuclear, eso llevó a lo satelital, lo satelital a los
radares''). ¿Qué significa trabajar por proyectos? ¿Por qué se la considera una empresa
estratégica?

\paragraph*{ITP Argentina}
Caso a desarrollar en clase con la información que aporte el docente: historia de la empresa,
qué fabrica o repara, sus clientes, sus certificaciones y su relación con el Área de Material
Quilmes y con la escuela.

\begin{recuadro}
\textbf{Guía para analizar un caso:} 1) origen y emprendedor; 2) problema u oportunidad que
detectó; 3) producto o servicio; 4) clientes y mercado; 5) recursos clave (técnicos,
humanos, financieros); 6) relación con el Estado y el territorio; 7) crisis y cómo las
enfrentó; 8) aprendizajes para nuestro proyecto.
\end{recuadro}

% ------------------------------------------------------------
\section*{3. EL SISTEMA CIENTÍFICO-TECNOLÓGICO ARGENTINO}

\subsubsection*{3.1. Jorge Sabato y la pregunta por la tecnología propia}
\textbf{Jorge Alberto Sabato} (1924--1983) fue un físico y tecnólogo argentino. Trabajó en la
\textbf{CNEA}, donde dirigió el área de metalurgia, impulsó en 1962 el \textbf{Servicio de
Asistencia Técnica a la Industria (SATI)} para transferir conocimiento a las empresas
nacionales y luego fue gerente de tecnología del organismo. El Instituto Sabato de la CNEA
lleva su nombre.

Su preocupación central era una pregunta que sigue vigente: \textbf{¿por qué los países de
América Latina compran casi toda la tecnología afuera, aunque tengan científicos bien
formados?} Su respuesta fue que no alcanza con tener buenos científicos, ni buenas empresas,
ni un Estado con recursos: el problema es que \textbf{esos tres actores no están conectados}.
Para explicarlo, en 1968 publicó junto con el politólogo \textbf{Natalio Botana} el trabajo
\textit{La ciencia y la tecnología en el desarrollo futuro de América Latina}, donde presentó
el modelo que hoy se conoce como \textbf{Triángulo de Sabato}.

\subsubsection*{3.2. Los tres vértices}

\begin{center}
\begin{tikzpicture}[scale=1.2, every node/.style={font=\small, align=center}]
  \coordinate (G) at (0,3.8);
  \coordinate (I) at (-3.8,0);
  \coordinate (E) at (3.8,0);
  \draw[very thick, ppAzul] (G) -- (I) -- (E) -- cycle;
  \fill[ppAzul] (G) circle (3.5pt) (I) circle (3.5pt) (E) circle (3.5pt);
  \node[above=3pt] at (G) {\textbf{G -- GOBIERNO}\\ (Estado)};
  \node[below=4pt] at (I) {\textbf{I -- INFRAESTRUCTURA}\\ \textbf{CIENTÍFICO-TECNOLÓGICA}};
  \node[below=4pt] at (E) {\textbf{E -- ESTRUCTURA}\\ \textbf{PRODUCTIVA}};
  \node[font=\footnotesize, rotate=45, anchor=south] at (-1.9,1.9) {financia y define prioridades};
  \node[font=\footnotesize, rotate=-45, anchor=south] at (1.9,1.9) {compra, regula y promueve};
  \node[font=\footnotesize, anchor=south] at (0,0.05) {conocimiento $\leftrightarrow$ demanda};
  \node[font=\footnotesize\itshape, ppAzul] at (0,1.4) {interrelaciones};
\end{tikzpicture}
\end{center}

\begin{center}
\small
\begin{tabular}{|p{0.22\linewidth}|p{0.38\linewidth}|p{0.30\linewidth}|}
\hline
\textbf{Vértice} & \textbf{Quiénes lo forman} & \textbf{Su rol} \\ \hline
\textbf{Gobierno (G)} & Poder Ejecutivo, ministerios, Fuerzas Armadas, agencias de
financiamiento, provincias y municipios. & Formula la \textbf{política} científica y
tecnológica, asigna recursos, regula y, muy importante, \textbf{compra} tecnología nacional. \\ \hline
\textbf{Infraestructura científico-tecnológica (I)} & Universidades, CONICET, CNEA, CONAE,
INTI, CITEDEF, laboratorios, institutos, escuelas técnicas. & \textbf{Genera conocimiento} y
forma a las personas (investigadores, ingenieros, técnicos). \\ \hline
\textbf{Estructura productiva (E)} & Empresas públicas y privadas, grandes y pymes, talleres,
cooperativas. & \textbf{Produce} bienes y servicios y \textbf{convierte el conocimiento en
productos} que se venden. \\ \hline
\end{tabular}
\end{center}

\subsubsection*{3.3. Las tres clases de relaciones}
Sabato decía que un vértice aislado no sirve: lo que importa son las \textbf{relaciones}.

\begin{enumerate}
  \item \textbf{Intrarrelaciones (dentro de cada vértice):} cada actor tiene que estar
        organizado por dentro. Por ejemplo, que las universidades y los laboratorios
        cooperen entre sí, o que las empresas de un sector se asocien en una cámara.
  \item \textbf{Interrelaciones (entre vértices, los lados del triángulo):}
  \begin{itemize}
    \item \textbf{G -- I:} el Estado financia la investigación y le marca prioridades
          (energía, defensa, salud, espacio).
    \item \textbf{G -- E:} el Estado actúa como \textbf{cliente}, regulador y promotor
          (compras públicas, créditos, beneficios impositivos, normas técnicas).
    \item \textbf{I -- E:} los laboratorios resuelven problemas de las empresas y las
          empresas les plantean demandas concretas; las universidades forman el personal que
          las empresas necesitan. \textbf{Es el lado más difícil de construir.}
  \end{itemize}
  \item \textbf{Extrarrelaciones (con el exterior):} cada vértice se vincula con el mundo:
        importar tecnología, exportar productos, cooperar con otros países, participar en
        organismos internacionales. Un triángulo fuerte importa tecnología \textbf{eligiendo}
        qué comprar y aprendiendo de ella, y también \textbf{exporta}.
\end{enumerate}

\subsubsection*{3.4. Cuando el triángulo no funciona}
\begin{center}
\small
\begin{tabular}{|p{0.36\linewidth}|p{0.54\linewidth}|}
\hline
\textbf{Síntoma} & \textbf{Qué falla} \\ \hline
La universidad investiga temas que ninguna empresa usa; los resultados quedan en papers. &
Lado \textbf{I -- E} roto: no hay transferencia de conocimiento. \\ \hline
Las empresas compran ``\textbf{paquetes tecnológicos}'' cerrados (máquinas, licencias,
software) y nunca aprenden a hacerlos ni a mejorarlos. & Extrarrelaciones dominantes:
dependencia tecnológica. \\ \hline
El Estado compra equipos importados aunque haya capacidad nacional para fabricarlos. & Lado
\textbf{G -- E} débil: el Estado no usa su poder de compra. \\ \hline
Proyectos que arrancan y se cortan con cada cambio de gobierno. & Falta de política de Estado
sostenida en el vértice \textbf{G}. \\ \hline
\end{tabular}
\end{center}

Sabato proponía \textbf{``abrir el paquete tecnológico''}: en lugar de comprar todo cerrado,
desarmar la tecnología en sus partes, comprar afuera solo lo que no se puede hacer y
desarrollar localmente el resto, aprendiendo en el proceso. Su mensaje central era que el
problema no es técnico sino \textbf{político}: alguien tiene que decidir construir el
triángulo y sostenerlo en el tiempo.

\subsubsection*{3.5. Ejemplos del triángulo funcionando}

\paragraph*{Radares argentinos.}
\textbf{G:} Ministerio de Defensa y Fuerza Aérea (definen requerimientos, financian y compran,
plan SINVICA). \textbf{I:} CNEA, Instituto Balseiro, universidades (ingenieros formados).
\textbf{E:} INVAP (diseña y fabrica). \textbf{Resultado:} el país pasó de importar radares a
fabricarlos, instalarlos en todo el territorio y exportarlos.

\paragraph*{Satélites.}
\textbf{G:} CONAE y ARSAT. \textbf{I:} universidades, centros de ensayo. \textbf{E:} INVAP y
pymes proveedoras. \textbf{Resultado:} satélites SAC-C, SAOCOM, ARSAT-1 y ARSAT-2.

\begin{recuadro}
\textbf{Caso internacional: el polo aeroespacial de São José dos Campos (Brasil).}
Es el ejemplo de manual del triángulo aplicado a la aviación. \textbf{G:} en 1946 la Fuerza
Aérea brasileña creó el Centro Técnico Aeroespacial (hoy DCTA) y el municipio cedió terrenos.
\textbf{I:} en 1950 se instaló allí el \textbf{ITA} (Instituto Tecnológico de Aeronáutica),
que formó a los ingenieros. \textbf{E:} en 1969, en ese mismo ambiente, nació
\textbf{Embraer}, hoy uno de los mayores fabricantes de aviones del mundo. A su alrededor se
formó un \textbf{clúster} de más de un centenar de empresas del sector. Fueron necesarios
décadas de política sostenida, formación de personas y compras del Estado.
\end{recuadro}

\paragraph*{Para pensar.} ¿Qué pasó con FAdeA? Tiene historia, personal capacitado y
capacidad técnica, pero en los últimos años perdió contratos y financiamiento. ¿Qué lado del
triángulo se debilitó?

\subsubsection*{3.6. Del triángulo a la triple hélice}
Desde los años noventa se habla de la \textbf{triple hélice}
(universidad -- empresa -- gobierno): la misma idea, con el énfasis en que los tres actores
se superponen y cada uno puede tomar funciones del otro (universidades que crean empresas,
empresas que investigan, Estado que invierte como socio). La \textbf{cuádruple hélice} agrega
a la \textbf{sociedad civil} (usuarios, comunidades, organizaciones sociales).

\subsubsection*{3.7. Los organismos del sistema}
\begin{center}
\small
\begin{tabular}{|p{0.15\linewidth}|p{0.75\linewidth}|}
\hline
\textbf{Organismo} & \textbf{Qué hace} \\ \hline
\textbf{CONICET} & Consejo Nacional de Investigaciones Científicas y Técnicas (1958). Principal
organismo de promoción de la ciencia: financia y emplea investigadores y becarios, y tiene
institutos en todo el país. \\ \hline
\textbf{CNEA} & Comisión Nacional de Energía Atómica (1950). Investigación y desarrollo nuclear:
reactores de investigación, combustibles, radioisótopos para medicina, centrales nucleares.
De ella dependen los Centros Atómicos (Constituyentes, Ezeiza, Bariloche) y los Institutos
Balseiro y Sabato; de ella nació INVAP. \\ \hline
\textbf{CONAE} & Comisión Nacional de Actividades Espaciales (1991). Agencia espacial argentina:
define y ejecuta el Plan Espacial Nacional (satélites de observación de la Tierra como SAC-C
y SAOCOM, y el desarrollo de lanzadores). INVAP es su principal contratista. \\ \hline
\textbf{CITEDEF} & Instituto de Investigaciones Científicas y Técnicas para la Defensa
(Villa Martelli). Investigación y desarrollo para la defensa y aplicaciones civiles: láseres,
sensores, propulsión, química, electrónica. \\ \hline
\textbf{INTI} & Instituto Nacional de Tecnología Industrial (1957). Asiste a la industria y a
las pymes: ensayos, \textbf{calibración} de instrumentos, metrología, certificación,
transferencia de tecnología y capacitación. \\ \hline
\end{tabular}
\end{center}

\paragraph*{¿Por qué le importa a un técnico?}
Estos organismos y las empresas que los rodean (INVAP, FAdeA, talleres, pymes proveedoras)
son \textbf{fuentes de empleo calificado}, de \textbf{prácticas profesionalizantes} y de
\textbf{servicios} que un emprendedor técnico necesita: el INTI calibra el instrumental, el
CONICET y las universidades resuelven problemas de investigación, y la CONAE y la Fuerza Aérea
son clientes de proyectos tecnológicos. Y la \textbf{escuela técnica} también es parte del
vértice I: forma a los técnicos que el triángulo necesita.

% ------------------------------------------------------------
\section*{4. DESARROLLO LOCAL Y TERRITORIO}

\subsubsection*{4.1. ¿Qué es el desarrollo local?}
El \textbf{desarrollo local} es un proceso de crecimiento económico y mejora de la calidad de
vida que se apoya en los recursos propios de un \textbf{territorio}: su gente, sus empresas,
sus escuelas, sus universidades, sus saberes y su historia. No espera que la solución venga
``de afuera'' (una gran inversión extranjera), sino que organiza lo que ya existe y lo conecta.
Es el Triángulo de Sabato aplicado a escala de un municipio o una región.

El \textbf{territorio} no es solo un lugar en el mapa: es un espacio con identidad, actores,
recursos e historia. Quilmes, por ejemplo, tiene una larga historia industrial (la cervecería,
la industria textil y metalúrgica, el propio predio de IMPA), una universidad nacional (UNQ),
escuelas técnicas, un parque industrial y un área de material de la Fuerza Aérea.

\subsubsection*{4.2. ¿Por qué las empresas se agrupan?}
Cuando muchas empresas se instalan cerca unas de otras obtienen ventajas que una empresa
aislada no tiene (los economistas las llaman \textbf{economías de aglomeración}):
\begin{itemize}
  \item \textbf{Infraestructura compartida:} energía, gas, agua, accesos para camiones,
        tratamiento de efluentes, seguridad. Más barata entre todos.
  \item \textbf{Proveedores cerca:} menos costos y tiempos de transporte.
  \item \textbf{Mano de obra calificada:} los técnicos se forman y se quedan en la zona.
  \item \textbf{Conocimiento que circula:} el personal cambia de empresa, se comparten
        prácticas, aparecen ideas nuevas.
  \item \textbf{Cooperación:} compras conjuntas, capacitación común, presencia en ferias,
        reclamos ante el Estado.
\end{itemize}
Según \textbf{qué} comparten y \textbf{cómo} se organizan, estas agrupaciones reciben nombres
distintos. Los más importantes son el \textbf{parque industrial}, el \textbf{polo
tecnológico} y el \textbf{clúster}. Muchas veces se confunden: veamos cada uno.

% ............................................................
\subsubsection*{4.3. El parque industrial}

\begin{recuadro}
\textbf{Definición:} un \textbf{parque industrial} es un \textbf{predio delimitado}, en una zona
habilitada para la industria, \textbf{loteado y dotado de infraestructura y servicios
comunes}, donde se instalan empresas manufactureras y de servicios. Lo administra un ente
(público, privado o mixto) y tiene un reglamento interno.
\end{recuadro}

\textbf{Una forma de imaginarlo:} es como un ``barrio planificado para fábricas''. Igual que un
barrio, tiene calles, lotes, servicios y reglas de convivencia, pero pensados para la
actividad industrial.

\paragraph*{Qué ofrece.}
\begin{itemize}
  \item \textbf{Lotes} o naves industriales para comprar o alquilar.
  \item \textbf{Infraestructura:} calles pavimentadas aptas para camiones, energía eléctrica
        (a veces en media tensión), gas, agua, desagües y tratamiento de efluentes,
        iluminación, cerco perimetral y seguridad.
  \item \textbf{Servicios comunes:} administración, control de acceso, balanza, estacionamiento,
        a veces comedor, salas de capacitación, bomberos, transporte para el personal.
  \item \textbf{Beneficios:} los municipios y la provincia suelen dar exenciones de tasas o
        impuestos y líneas de crédito a quienes se radican en un agrupamiento habilitado.
  \item \textbf{Orden urbano y ambiental:} separa las industrias de las viviendas y facilita el
        control ambiental.
\end{itemize}

\paragraph*{Marco legal en la Provincia de Buenos Aires.}
La \textbf{Ley 13.744} regula los \textbf{agrupamientos industriales} y distingue varias
categorías:
\begin{itemize}
  \item \textbf{Parque industrial:} con toda la infraestructura y los servicios completos.
  \item \textbf{Sector industrial planificado:} una porción de zona industrial con
        infraestructura mínima; puede transformarse en parque cuando completa las obras.
  \item \textbf{Área de servicios industriales y logística:} para depósitos, transporte y
        distribución.
  \item \textbf{Incubadoras de empresas}, \textbf{unidades modulares productivas} y
        \textbf{parques industriales pequeños y medianos} (agregados por la Ley 14.792).
\end{itemize}
Según quién los impulsa pueden ser \textbf{oficiales} (Nación, provincia o municipio),
\textbf{privados} o \textbf{mixtos}.

\paragraph*{Ejemplos cercanos.}
\begin{itemize}
  \item \textbf{Parque Industrial La Bernalesa (Quilmes):} recuperó un predio industrial
        histórico; allí se desarrolla el programa \textbf{Quilmes Tech} de la Universidad
        Nacional de Quilmes y el municipio, y se proyecta un polo tecnológico.
  \item \textbf{Parque Industrial y Tecnológico de Florencio Varela (PITEC).}
  \item \textbf{Parque Industrial Pilar}, uno de los más grandes del país, con predominio de
        grandes empresas.
\end{itemize}
En la provincia hay alrededor de \textbf{200 agrupamientos industriales}, repartidos en más de
100 municipios.

\paragraph*{Lo importante:} en un parque industrial las empresas pueden ser de
\textbf{rubros totalmente distintos} (una fábrica de envases al lado de una metalúrgica y un
laboratorio). Lo que comparten es el \textbf{espacio y la infraestructura}, no
necesariamente el negocio.

% ............................................................
\subsubsection*{4.4. El polo (o parque) tecnológico}

\begin{recuadro}
\textbf{Definición:} un \textbf{polo tecnológico} es un espacio donde conviven, a propósito,
\textbf{instituciones de conocimiento} (universidades, centros de investigación, laboratorios
públicos) y \textbf{empresas de base tecnológica}, con el objetivo de que el conocimiento se
transforme en productos, servicios y nuevas empresas.
\end{recuadro}

\textbf{Una forma de imaginarlo:} es como un ``campus'' donde la facultad, el laboratorio y la
empresa están puerta con puerta. Es el Triángulo de Sabato hecho edificio.

\paragraph*{Qué lo diferencia de un parque industrial.}
\begin{itemize}
  \item Lo que se busca no es fabricar en serie, sino \textbf{investigar, desarrollar,
        innovar} y crear empresas nuevas.
  \item Siempre hay una o más \textbf{instituciones científicas} como ``ancla''.
  \item Suele incluir \textbf{incubadoras}, laboratorios compartidos, salas limpias, centros de
        ensayo y oficinas de transferencia de tecnología.
  \item El recurso principal no es la nave industrial sino el \textbf{conocimiento} y las
        \textbf{personas} (investigadores, ingenieros, técnicos).
\end{itemize}

\paragraph*{Ejemplos.}
\begin{itemize}
  \item \textbf{Polo Tecnológico Constituyentes} (San Martín, provincia de Buenos Aires): reúne
        organismos como el Centro Atómico Constituyentes de la CNEA, el INTI y CITEDEF, y
        universidades de la zona.
  \item \textbf{Bariloche:} alrededor del Centro Atómico, el Instituto Balseiro e INVAP se
        formó uno de los principales polos científicos y tecnológicos del país.
  \item \textbf{Quilmes:} el proyecto de polo tecnológico en La Bernalesa, junto a la UNQ.
\end{itemize}

% ............................................................
\subsubsection*{4.5. El clúster}

\begin{recuadro}
\textbf{Definición:} un \textbf{clúster} (del inglés \textit{cluster}, ``racimo'') es una
\textbf{concentración geográfica de empresas de un mismo sector y de sus actividades
relacionadas} (proveedores, clientes, prestadores de servicios, instituciones de formación e
investigación, cámaras) que \textbf{compiten y cooperan al mismo tiempo}. El concepto fue
difundido por el economista Michael Porter.
\end{recuadro}

\textbf{Una forma de imaginarlo:} no es un predio con cerco sino una \textbf{red}. Es como un
racimo de uvas: muchas empresas distintas, unidas por el mismo ``tallo'' (el sector y la
región). Pueden estar dispersas en toda una ciudad o una región.

\paragraph*{Cómo funciona.}
\begin{itemize}
  \item Suele haber una o varias \textbf{empresas ancla} (una gran fábrica) rodeadas de
        \textbf{pymes proveedoras} especializadas.
  \item Las empresas \textbf{compiten} por clientes, pero \textbf{cooperan} en lo que les
        conviene a todas: formar técnicos, certificar calidad, comprar insumos en conjunto,
        exportar juntas, pedir infraestructura al Estado.
  \item Se forma un \textbf{mercado de trabajo especializado}: los técnicos de la zona saben
        del sector y circulan entre empresas.
  \item Cuando se organiza formalmente (con una asociación, una mesa de trabajo, una
        gerencia) se habla de \textbf{clúster organizado} o, en Brasil, de
        \textbf{Arranjo Produtivo Local}.
\end{itemize}

\paragraph*{Ejemplos.}
\begin{itemize}
  \item \textbf{Clúster aeroespacial de São José dos Campos (Brasil):} Embraer como ancla, el
        ITA y el DCTA como instituciones de conocimiento y más de cien empresas del sector.
  \item \textbf{Clúster aeronáutico de Córdoba:} FAdeA y las pymes metalmecánicas y de
        servicios que trabajan para ella. Su crisis muestra el riesgo de depender de una sola
        empresa ancla.
  \item \textbf{Toulouse (Francia):} alrededor de Airbus, con miles de proveedores, escuelas de
        ingeniería y centros de investigación.
  \item \textbf{Silicon Valley (EE.\,UU.):} el clúster tecnológico más conocido del mundo.
  \item Fuera de la tecnología: el clúster del vino en Mendoza, el del mueble, el del calzado o
        el del software en distintas ciudades argentinas.
\end{itemize}

% ............................................................
\subsubsection*{4.6. Cuadro comparativo}
\begin{center}
\footnotesize
\setlength{\tabcolsep}{3pt}
\begin{tabular}{|>{\raggedright\arraybackslash}p{0.17\linewidth}|>{\raggedright\arraybackslash}p{0.25\linewidth}|>{\raggedright\arraybackslash}p{0.25\linewidth}|>{\raggedright\arraybackslash}p{0.25\linewidth}|}
\hline
 & \textbf{Parque industrial} & \textbf{Polo tecnológico} & \textbf{Clúster} \\ \hline
\textbf{¿Qué es?} & Un \textbf{predio} con infraestructura. & Un \textbf{espacio} de
conocimiento e innovación. & Una \textbf{red} de empresas e instituciones. \\ \hline
\textbf{¿Tiene límites físicos?} & Sí, cerco y lotes. & Generalmente sí (campus, edificio,
predio). & No: abarca una ciudad o región. \\ \hline
\textbf{¿Qué comparten?} & Infraestructura y servicios. & Conocimiento, laboratorios,
investigadores. & Sector, proveedores, mano de obra, cooperación. \\ \hline
\textbf{¿Mismo rubro?} & No necesariamente. & Base tecnológica, rubros variados. & Sí, un sector
y sus actividades relacionadas. \\ \hline
\textbf{¿Quién lo organiza?} & Ente administrador (público, privado o mixto). & Universidades,
organismos científicos, Estado. & Muchas veces surge solo; puede organizarse con una
asociación. \\ \hline
\textbf{Actividad principal} & Producción, logística. & Investigación, desarrollo,
incubación. & Toda la cadena de valor del sector. \\ \hline
\textbf{Ejemplo} & La Bernalesa (Quilmes), Pilar. & Constituyentes, Bariloche. & São José dos
Campos, Córdoba aeronáutico. \\ \hline
\end{tabular}
\end{center}

\paragraph*{Pueden superponerse.} Un clúster puede tener empresas dentro de un parque
industrial y otras fuera; un parque industrial puede incluir un polo tecnológico (como se
proyecta en La Bernalesa); y un polo tecnológico puede ser el ``cerebro'' de un clúster (como
el ITA en São José dos Campos).

% ............................................................
\subsubsection*{4.7. Otros conceptos del territorio}
\begin{itemize}
  \item \textbf{Zona industrial:} área que el código de planeamiento municipal habilita para
        uso industrial. Es la base: los agrupamientos se ubican dentro de ellas.
  \item \textbf{Cadena de valor:} todas las etapas por las que pasa un producto, desde la
        materia prima hasta el cliente final, con las empresas que participan en cada una.
        Ejemplo aeronáutico: fabricante de componentes $\rightarrow$ integrador de aviónica
        $\rightarrow$ fabricante de aeronaves $\rightarrow$ operador aéreo $\rightarrow$
        taller de mantenimiento $\rightarrow$ desguace y reciclado. Un clúster fuerte reúne
        varias etapas de la cadena en la misma región.
  \item \textbf{Empresa ancla:} una empresa grande que, por su demanda, sostiene a muchas
        pymes proveedoras (FAdeA, Embraer, Airbus, INVAP).
  \item \textbf{Cámaras sectoriales:} agrupan a empresas de un rubro para representarlas y
        cooperar, como la Cámara Argentina Aeronáutica y Espacial (CARAE).
\end{itemize}

\subsubsection*{4.8. Herramientas de promoción del desarrollo local}
\begin{itemize}
  \item \textbf{Planificación estratégica participativa:} el municipio, las empresas, las
        escuelas y las organizaciones sociales definen juntos un diagnóstico y prioridades
        para el territorio.
  \item \textbf{Agencias de desarrollo:} organismos (municipales, regionales o mixtos) que
        asesoran a emprendedores, gestionan créditos y programas y vinculan empresas.
  \item \textbf{Creación de agrupamientos industriales y polos tecnológicos:} el Estado aporta
        tierra, obras de infraestructura y beneficios para atraer empresas.
  \item \textbf{Incubadoras de empresas:} acompañan a emprendimientos nuevos con espacio
        físico, capacitación, mentores y contactos durante sus primeros años. Pueden ser
        \textbf{tecnológicas} (universidades, polos), \textbf{sociales} (economía popular,
        cooperativas) o \textbf{culturales}. Las \textbf{aceleradoras} trabajan con proyectos
        más avanzados y suelen invertir a cambio de una participación.
  \item \textbf{Microemprendimientos y emprendimientos familiares:} con apoyo en capacitación,
        microcréditos y ferias.
  \item \textbf{Asociativismo intermunicipal y microrregiones:} varios municipios cercanos
        que se unen para encarar proyectos que solos no podrían (un parque industrial
        regional, una planta de reciclado, un aeródromo).
  \item \textbf{Vinculación escuela--empresa:} Prácticas Profesionalizantes, pasantías,
        proyectos para empresas locales. La escuela técnica es un actor del desarrollo local.
\end{itemize}

\subsubsection*{4.9. Desarrollo rural y sustentabilidad}
El desarrollo no es solo urbano e industrial: el campo genera trabajo y necesita tecnología
(aviación agrícola, drones de pulverización y monitoreo, electrónica para maquinaria,
conectividad rural). Un desarrollo \textbf{sustentable} busca crecer sin agotar los recursos
ni dañar el ambiente, cuidando el agua, el suelo y el aire para las generaciones futuras.
% ------------------------------------------------------------
\section*{5. ACTIVIDADES DE LA UNIDAD}

\begin{enumerate}
  \item \textbf{Análisis de casos:} en grupos, analizar uno de los casos aeronáuticos con la
        guía del recuadro y exponerlo al curso.
  \item \textbf{Mapa productivo de Quilmes y alrededores:} en un mapa (papel o Google My
        Maps) ubicar empresas, talleres, parques industriales, escuelas técnicas,
        universidades y organismos de la zona. Clasificar cada lugar: ¿es parte de un parque
        industrial, de un polo tecnológico, de un clúster? ¿Qué cadenas de valor existen?
        ¿Dónde podría insertarse un técnico aviónico?
  \item \textbf{¿Parque, polo o clúster?} Para cada caso, decidir de qué se trata y
        justificar: a) un predio con 40 lotes, calles, energía y seguridad, con fábricas de
        distintos rubros; b) un edificio donde una universidad comparte laboratorios con
        empresas de software y electrónica; c) más de cien empresas aeronáuticas dispersas en
        una región, que comparten proveedores y forman técnicos en conjunto; d) La Bernalesa
        con su proyecto Quilmes Tech.
  \item \textbf{Diseñar un parque o polo:} en grupos, proponer al municipio un agrupamiento
        para el sector aeronáutico y electrónico del sur del conurbano: ubicación,
        infraestructura, servicios comunes, instituciones ancla y beneficios para atraer
        empresas.
  \item \textbf{Entrevista a un emprendedor:} entrevistar a un emprendedor del barrio o del
        sector y relacionar su historia con las teorías y los tipos de emprendimiento.
  \item \textbf{Triángulo de Sabato:} elegir un desarrollo tecnológico argentino (un radar,
        un satélite, un reactor, una vacuna) y dibujar su triángulo: quién fue el Estado, quién
        la infraestructura científica y quién la estructura productiva. ¿Qué lado del triángulo
        falló o funcionó?
  \item \textbf{Visita técnica:} recorrer un organismo o empresa del Área de Material Quilmes o
        del Polo Tecnológico Constituyentes e informar qué hace, qué perfiles técnicos emplea y
        cómo se vincula con otros actores.
  \item \textbf{Debate:} ``¿Se nace o se hace emprendedor?'' y ``¿Emprender es siempre una
        elección?''.
\end{enumerate}
